% GENERATED FILE. Edit canonical lessons, then run npm run generate:books.
% canonical-source-hash: fnv1a64:168dcbd37e6b3290
% canonical-lessons: TA-C380-poruppu, TA-C380-katappatu, TA-C380-arppani, TA-C380-pankali, TA-C380-ottulai

\chapter{Duty and Commitment}
\label{ch:ta-duty-and-commitment}
\hlchaptermodality{380}

\begin{chapteropening}
\textbf{By the end of this chapter:} \emph{I can talk about responsibility, obligation and contribution.}

Complete the last lesson of chapter 380: I can talk about responsibility, obligation and contribution.
\end{chapteropening}

\section[\texorpdfstring{\ta{பொறுப்பு} (poṟuppu)}{பொறுப்பு (poṟuppu)}]{\ta{பொறுப்பு} (poṟuppu) — responsibility}

\label{lesson:TA-C380-poruppu}



\begin{quote}
Before the new one: say the Tamil for to muster, then the Tamil for welfare.
\end{quote}

\subsection*{You'll want to know: \ta{பொறுப்பு}}
\textbf{\ta{பொறுப்பு}} — \emph{poṟuppu} — \textquotedblleft{}responsibility\textquotedblright{}.

\begin{quote}
\ta{குழந்தைகளைப்} \ta{பார்த்துக்கொள்வது} \ta{பெற்றோரின்} \ta{பொறுப்பு}.
\end{quote}

\begin{quote}
\emph{kuḻantaikaḷaip pārttukkoḷvatu peṟṟōriṉ poṟuppu.}
\end{quote}

\begin{quote}
Looking after children is the parents' responsibility.
\end{quote}

\ta{பொறுப்பாக}, responsibly. \ta{பொறுப்பேற்றுக்கொள்}, to take responsibility.

\begin{cousinweb}
\textbf{\ta{பொறுப்பு}} (\emph{poṟuppu}) is native Dravidian: DEDR 4565 lists Tamil \emph{poṟu}, \textquotedblleft{}to bear, sustain, endure, take responsibility\textquotedblright{}, and \emph{poṟuppu}, \textquotedblleft{}a burden borne\textquotedblright{}.
\end{cousinweb}

\begin{grammarlens}[title={What you've built}]
One more word for giving reasons and opinions.
\end{grammarlens}

\subsection*{Your turn}
\begin{itemize}
\raggedright
  \item \emph{Say it:} \emph{poṟuppu}
  \item \emph{Say it:} \emph{poṟuppu}, once more
  \item \emph{Say it:} \emph{kuḻantaikaḷaip pārttukkoḷvatu peṟṟōriṉ poṟuppu.} — Looking after children is the parents' responsibility.
  \item \emph{Recall:} say the Tamil for to muster, then the Tamil for welfare, then say \emph{poṟuppu} again
\end{itemize}

\begin{tcolorbox}[breakable,colback=teal!4,colframe=teal!35!black,title={Before you move on}]
What does \ta{பொறுப்பு} mean? (\textquotedblleft{}Responsibility\textquotedblright{}.) Say it once more.
\end{tcolorbox}
\section[\texorpdfstring{\ta{கடப்பாடு} (kaṭappāṭu)}{கடப்பாடு (kaṭappāṭu)}]{\ta{கடப்பாடு} (kaṭappāṭu) — an obligation; a commitment}

\label{lesson:TA-C380-katappatu}



\begin{quote}
Before the new one: say the Tamil for welfare, then the Tamil for responsibility.
\end{quote}

\subsection*{You'll want to know: \ta{கடப்பாடு}}
\textbf{\ta{கடப்பாடு}} — \emph{kaṭappāṭu} — \textquotedblleft{}an obligation; a commitment\textquotedblright{}.

\begin{quote}
\ta{கொடுத்த} \ta{வாக்கைக்} \ta{காப்பாற்றுவது} \ta{நம்} \ta{கடப்பாடு}.
\end{quote}

\begin{quote}
\emph{koṭutta vākkaik kāppāṟṟuvatu nam kaṭappāṭu.}
\end{quote}

\begin{quote}
Keeping our word is our obligation.
\end{quote}

A formal word close to \ta{கடமை}, duty; \ta{கடமை} and \ta{கடன்}, a debt, share one root.

\begin{cousinweb}
\textbf{\ta{கடப்பாடு}} (\emph{kaṭappāṭu}) is native: DEDR 1113 lists Tamil \emph{kaṭaṉ}, \textquotedblleft{}obligation, debt\textquotedblright{}, and \emph{kaṭamai}, \textquotedblleft{}duty\textquotedblright{}; the word is usually read as that root joined to \emph{pāṭu}, the noun of \emph{paṭu}: the state of being bound. DEDR does not list \emph{kaṭappāṭu} itself.
\end{cousinweb}

\begin{grammarlens}[title={What you've built}]
One more word for giving reasons and opinions.
\end{grammarlens}

\subsection*{Your turn}
\begin{itemize}
\raggedright
  \item \emph{Say it:} \emph{kaṭappāṭu}
  \item \emph{Say it:} \emph{kaṭappāṭu}, once more
  \item \emph{Say it:} \emph{koṭutta vākkaik kāppāṟṟuvatu nam kaṭappāṭu.} — Keeping our word is our obligation.
  \item \emph{Recall:} say the Tamil for welfare, then the Tamil for responsibility, then say \emph{kaṭappāṭu} again
\end{itemize}

\begin{tcolorbox}[breakable,colback=teal!4,colframe=teal!35!black,title={Before you move on}]
What does \ta{கடப்பாடு} mean? (\textquotedblleft{}An obligation; a commitment\textquotedblright{}.) Say it once more.
\end{tcolorbox}
\section[\texorpdfstring{\ta{அர்ப்பணி} (arppaṇi)}{அர்ப்பணி (arppaṇi)}]{\ta{அர்ப்பணி} (arppaṇi) — to dedicate; to devote}

\label{lesson:TA-C380-arppani}



\begin{quote}
Before the new one: say the Tamil for responsibility, then the Tamil for an obligation.
\end{quote}

\subsection*{You'll want to know: \ta{அர்ப்பணி}}
\textbf{\ta{அர்ப்பணி}} — \emph{arppaṇi} — \textquotedblleft{}to dedicate; to devote\textquotedblright{}.

\begin{quote}
\ta{அவர்} \ta{தன்} \ta{வாழ்க்கையைக்} \ta{கல்விக்காக} \ta{அர்ப்பணித்தார்}.
\end{quote}

\begin{quote}
\emph{avar taṉ vāḻkkaiyaik kalvikkāka arppaṇittār.}
\end{quote}

\begin{quote}
He devoted his life to education.
\end{quote}

Past (I …): \ta{அர்ப்பணித்தேன்}, \emph{arppaṇittēṉ}.

What is given takes -\ta{ஐ}, the cause -\ta{க்கு} or -\ta{க்காக}.

\begin{cousinweb}
\textbf{\ta{அர்ப்பணி}} (\emph{arppaṇi}) is a verb built on Sanskrit \emph{arpaṇa}, \textquotedblleft{}fixing; placing upon\textquotedblright{} (CDIAL 683a), whose Prakrit form means \textquotedblleft{}gift, offering\textquotedblright{}.
\end{cousinweb}

\begin{grammarlens}[title={What you've built}]
One more word for giving reasons and opinions.
\end{grammarlens}

\subsection*{Your turn}
\begin{itemize}
\raggedright
  \item \emph{Say it:} \emph{arppaṇi}
  \item \emph{Say it:} \emph{arppaṇi}, once more
  \item \emph{Say it:} \emph{avar taṉ vāḻkkaiyaik kalvikkāka arppaṇittār.} — He devoted his life to education.
  \item \emph{Recall:} say the Tamil for responsibility, then the Tamil for an obligation, then say \emph{arppaṇi} again
\end{itemize}

\begin{tcolorbox}[breakable,colback=teal!4,colframe=teal!35!black,title={Before you move on}]
What does \ta{அர்ப்பணி} mean? (\textquotedblleft{}To dedicate; to devote\textquotedblright{}.) Say it once more.
\end{tcolorbox}
\section[\texorpdfstring{\ta{பங்களி} (paṅkaḷi)}{பங்களி (paṅkaḷi)}]{\ta{பங்களி} (paṅkaḷi) — to contribute}

\label{lesson:TA-C380-pankali}



\begin{quote}
Before the new one: say the Tamil for an obligation, then the Tamil for to dedicate.
\end{quote}

\subsection*{You'll want to know: \ta{பங்களி}}
\textbf{\ta{பங்களி}} — \emph{paṅkaḷi} — \textquotedblleft{}to contribute\textquotedblright{}.

\begin{quote}
\ta{ஒவ்வொருவரும்} \ta{இந்த} \ta{வெற்றிக்குப்} \ta{பங்களித்தார்கள்}.
\end{quote}

\begin{quote}
\emph{ovvoruvarum inta veṟṟikkup paṅkaḷittārkaḷ.}
\end{quote}

\begin{quote}
Everyone contributed to this success.
\end{quote}

Past (I …): \ta{பங்களித்தேன்}, \emph{paṅkaḷittēṉ}.

To what: -\ta{க்கு}. The noun is \ta{பங்களிப்பு}, a contribution.

\begin{cousinweb}
\textbf{\ta{பங்களி}} (\emph{paṅkaḷi}) is native: \emph{paṅku}, \textquotedblleft{}share, portion\textquotedblright{} (DEDR 3808), and \emph{aḷi}, \textquotedblleft{}to show favour, to grant\textquotedblright{} (DEDR 301, where Malayalam \emph{aḷikkuka} means \textquotedblleft{}to bestow\textquotedblright{}): to grant one's share.
\end{cousinweb}

\begin{grammarlens}[title={What you've built}]
One more word for giving reasons and opinions.
\end{grammarlens}

\subsection*{Your turn}
\begin{itemize}
\raggedright
  \item \emph{Say it:} \emph{paṅkaḷi}
  \item \emph{Say it:} \emph{paṅkaḷi}, once more
  \item \emph{Say it:} \emph{ovvoruvarum inta veṟṟikkup paṅkaḷittārkaḷ.} — Everyone contributed to this success.
  \item \emph{Recall:} say the Tamil for an obligation, then the Tamil for to dedicate, then say \emph{paṅkaḷi} again
\end{itemize}

\begin{tcolorbox}[breakable,colback=teal!4,colframe=teal!35!black,title={Before you move on}]
What does \ta{பங்களி} mean? (\textquotedblleft{}To contribute\textquotedblright{}.) Say it once more.
\end{tcolorbox}
\section[\texorpdfstring{\ta{ஒத்துழை} (ottuḻai)}{ஒத்துழை (ottuḻai)}]{\ta{ஒத்துழை} (ottuḻai) — to cooperate}

\label{lesson:TA-C380-ottulai}



\begin{quote}
Before the new one: say the Tamil for to dedicate, then the Tamil for to contribute.
\end{quote}

\subsection*{You'll want to know: \ta{ஒத்துழை}}
\textbf{\ta{ஒத்துழை}} — \emph{ottuḻai} — \textquotedblleft{}to cooperate\textquotedblright{}.

\begin{quote}
\ta{இரண்டு} \ta{குழுக்களும்} \ta{ஒத்துழைத்து} \ta{வேலை} \ta{செய்தன}.
\end{quote}

\begin{quote}
\emph{iraṇṭu kuḻukkaḷum ottuḻaittu vēlai ceytaṉa.}
\end{quote}

\begin{quote}
The two teams worked together.
\end{quote}

Past (I …): \ta{ஒத்துழைத்தேன்}, \emph{ottuḻaittēṉ}.

With -\ta{உடன்}: \ta{அவருடன்} \ta{ஒத்துழை}. The noun is \ta{ஒத்துழைப்பு}, cooperation.

\begin{cousinweb}
\textbf{\ta{ஒத்துழை}} (\emph{ottuḻai}) is native: \emph{ottu}, \textquotedblleft{}agreeing, suiting\textquotedblright{}, from \emph{o} (DEDR 924), the verb inside \ta{ஒத்துக்கொள்}, and \emph{uḻai}, \textquotedblleft{}to labour hard\textquotedblright{} (DEDR 680): to work in agreement.
\end{cousinweb}

\begin{grammarlens}[title={What you've built}]
One more word for giving reasons and opinions.
\end{grammarlens}

\subsection*{Your turn}
\begin{itemize}
\raggedright
  \item \emph{Say it:} \emph{ottuḻai}
  \item \emph{Say it:} \emph{ottuḻai}, once more
  \item \emph{Say it:} \emph{iraṇṭu kuḻukkaḷum ottuḻaittu vēlai ceytaṉa.} — The two teams worked together.
  \item \emph{Recall:} say the Tamil for to dedicate, then the Tamil for to contribute, then say \emph{ottuḻai} again
\end{itemize}

\begin{tcolorbox}[breakable,colback=teal!4,colframe=teal!35!black,title={Before you move on}]
What does \ta{ஒத்துழை} mean? (\textquotedblleft{}To cooperate\textquotedblright{}.) Say it once more.
\end{tcolorbox}
